\textbf{Instalación y decisión eléctrica.} Se proyecta una instalación nueva de 400/230 V, 50 Hz, tres fases, neutro y PE, con generación, transferencia, UPS, banco y distribución propios en A y B. Por camino se especifica Eaton \texttt{93T60KMBSB}, de 60 kVA/60 kW, sin baterías internas, y FG Wilson \texttt{P200-6}, alternador estándar \texttt{FGL30120}, de 180 kVA/144 kW PRP. El CLIENTE adquiere hardware y ejecuta las obras asignadas; Synaptix especifica, coordina, integra y administra. La tensión y las dimensiones corresponden al proyecto, sin atribuirlas al tablero existente \parencites[RT-06.06--13, p.~15]{pucv-transversal}{p7-eaton-93t60}{p7-fg-p200-ficha}.

La selección de 60 kW incorpora el parque completo y la envolvente de una unidad climática mayor. La combinación climática de referencia es \texttt{s-MEXT-G00-DX-F2-022-S + PUZ-ZM250YKA2}; conserva condición de candidata hasta completar temperatura, humedad y calor por estado. No se autoriza su compra desde el punto nominal. Las UPS de 40 kVA y la unidad 013 resultan insuficientes para esta envolvente eléctrica o para la cota térmica con recalentamiento; no constituyen la compra propuesta. Las alternativas monofásicas 9PX y los grupos de 6 kVA tampoco cubren la carga completa \parencites{p6-eaton-tecnica40}{p6-mext-pi-2025}{p7-eaton-93t-tecnica}.

\textbf{Frontera completa de adquisición.} El suministro protege dieciséis PC, treinta y dos monitores, cómputo, red, CPE, conversión, medición, accesos, radio/CCTV, carga móvil y auxiliares. Los siete puestos de prioridad de negocio no excluyen los otros nueve. La Tabla~\ref{tab:4-balance-parcial} resume cotas de adquisición e integración; T-11 las descompone por familia y estado. Una BOM que las exceda exige sustitución o recálculo antes de compra. El veinte por ciento se aplica después de completar las partidas; no representa un consumo de catálogo ni absorbe una carga desconocida.

\begin{table}[htbp]
\centering\small
\caption{Envolvente eléctrica proyectada y crecimiento}\label{tab:4-balance-parcial}
\begin{tabularx}{\linewidth}{L{\dimexpr.40\linewidth-2\tabcolsep\relax}R{\dimexpr.20\linewidth-2\tabcolsep\relax}R{\dimexpr.20\linewidth-2\tabcolsep\relax}X}
\hline
\EncabezadoTabla{Familia} & \EncabezadoTabla{Inicial kW} & \EncabezadoTabla{Crecimiento kW} & \EncabezadoTabla{Frontera} \\ \hline
Central, protección y radio & 3,0100 & 3,0100 & Entrada AC \\ \hline
16 PC, 32 monitores y periféricos & 5,9200 & 5,9200 & Entrada AC \\ \hline
DC/PoE/medición convertido & 4,2142 & 5,0342 & PD y pérdidas una vez \\ \hline
Borde, motores y carga móvil & 2,6000 & 2,6000 & Entrada AC \\ \hline
Total sin climatización & 15,7442 & 16,5642 & Parque completo \\ \hline
Envolvente de un circuito 022 & 23,8455 & 23,8455 & Cota activa $P\leq S$ \\ \hline
Camino superviviente con 20\,\% & 47,5076 & 48,4916 & Salida de diseño \\ \hline
\end{tabularx}
\FuenteTabla{Cálculo de Synaptix desde T-11. El circuito climático deriva de $230(35,9)/1000+\sqrt{3}(400)(22,5)/1000=23,845457$ kVA \parencite[p.~3]{p6-mext-pi-2025}; la cota activa no atribuye FP unitario.}
\end{table}

La medición crece de 320 a 380 medidores de cada servicio; CCTV, de treinta y ocho a sesenta cámaras; el piloto conserva cuarenta sensores. DC/PoE suma 3,7028/4,4408 kW, incluyendo cable, y se convierte una vez con eficiencia mínima de compra 0,90 y 0,100 kW adicionales de espera. Los veinte vatios auxiliares por puesto corresponden a periféricos externos; teclado/mouse USB incluidos en el PC no se suman otra vez. El NOC remoto, Travelift, bombas industriales, embarcaciones y restaurante mantienen balances separados.

La capacidad UPS de 60 kVA se contrasta con 51,2335 kVA de diseño en crecimiento, que ya incluyen el veinte por ciento: la reserva adicional de placa es $1-51,2335/60=14,61\,\%$. El banco se dimensiona a plena carga y el grupo desde la envolvente del rectificador, no desde el promedio anual. El P200-6 de 144 kW nominales deja $1-129/144=10,42\,\%$ hasta el mínimo efectivo exigido de 129 kW; ese margen no demuestra derating admisible. La baja demanda anual exige comprobar el régimen de baja carga y programar los ensayos externos del grupo según fabricante antes de habilitación. La decisión conserva capacidad de falla y recuperación, y deberá revisarse si la BOM o la selección climática reducen la envolvente; no se considera libre de revisión por cumplir un techo nominal \parencite[RT-15.01, p.~27]{pucv-transversal}.

\textbf{Potencia aparente y reparto por fase.} Se fijan FP mínimos de compra 0,90 para central, puestos, conversión y borde, y 0,75 para motores y carga móvil. La suma no climática de crecimiento es 18,849136 kVA. El clima ya está expresado en kVA y no se vuelve a dividir por FP:
\[
P_d=1,20(16,564222+23,845457)=48,491615\ \mathrm{kW},\qquad
S_d=1,20(18,849136+23,845457)=51,233512\ \mathrm{kVA}.
\]
El cociente de diseño $P_d/S_d=0,9465$ procede de sumas conservadoras de magnitudes, no de una medición fasorial. El interior 022 se alimenta en L1 y el exterior en las tres fases, desde el mismo camino; se conservan circuitos, neutros y comunicación según su manual. Los dieciséis puestos se reparten 0/8/8 y los juegos móviles 6/5/5; las cargas centrales pasan a 0/810/2.200 W por fase. T-11 conserva cada reparto PD, conversión y cable.

La demanda nominal de crecimiento es 19,3965/14,4140/17,4229 kVA por fase, inferior a 20 kVA por fase. Se usa 230 V como referencia conservadora frente a $400/\sqrt{3}$; con $-2\,\%$, 225,4 V, las máximas de régimen, bancada o arranque son 84,511/66,479/80,893 A, inferiores a los 86 A publicados. Con $+2\,\%$, corrientes climáticas constantes y no-clima como potencia constante, los máximos permanecen bajo 20 kVA/kW por fase. Son cotas estacionarias de selección; la recepción captura cresta, escalones, desbalance y corriente real de las variantes \parencites[pp.~3--4]{p7-eaton-93t-tecnica}[pp.~33--35]{p7-eaton-93t-guia}.

Cada camino superviviente sostiene todo el parque no climático y su propio circuito de clima. El relevo de los dos climatizadores ocurre en UPS distintas. No se fuerza reparto normal 50/50 de PSU. La bancada de reservas se limita a 1,0 kW/1,25 kVA, con 0/0,60/0,65 kVA por fase, y se inhibe durante arranque de barrera. La otra barrera debe estar en reposo antes de iniciar, sin detener una maniobra en curso ni retirar lectores, seguridad o liberación de emergencia. Se exige entrada agregada de ambos circuitos de motor $\leq2,4$ kVA durante arranque, incluida espera no contada en control DC, frente a 1,6 kVA de régimen; el enclavamiento y estado seguro se verifican en las controladoras.

\textbf{Modo regulado y bypass.} La guía 93T documenta \textit{Frequency Converter} como operación equivalente a Online/batería con bypass no disponible. Se especifica entrada y salida de 50 Hz, ESS inhabilitado, Auto Bypass de arranque deshabilitado y alimentación separada de bypass con seccionador abierto y bloqueado en servicio. Se confirma inversor antes de habilitar cargas; no se atribuye al bypass la regulación del inversor ni una ventana ajustable no documentada. Mantenimiento exige camino sano verificado, parque transferido y clima propio detenido. La pérdida del inversor conserva servicio por la otra fuente/camino de las cargas que la acrediten; un activo de entrada única mantiene su brecha individual y no se convierte en doble PSU mediante un STS \parencite[pp.~21/27/80]{p7-eaton-93t-guia}.

\textbf{Banco y autonomía mantenida.} Cada UPS utiliza dos cadenas externas homogéneas de cuarenta CSB \texttt{HRL12540W}, revisión \texttt{RA260131}, VRLA-AGM de 12 V, con neutro DC al centro 20/20, positivo y negativo. Son ochenta bloques por UPS y ciento sesenta en A/B. La guía admite cantidad par de 32--44 bloques y recomienda protección \texttt{NZMN2-4-A200}, referencia \texttt{265874}, para 60 kW; la interfaz de disparo TB20 es de 24 V DC. El ensamble incorpora protección, supervisión y seccionamiento de cada cadena, coordinación y cableado equipotencial. No se atribuye compatibilidad al gabinete EBC-E de dos terminales \parencite[pp.~35/49--51]{p7-eaton-93t-guia}.

CSB publica 1.129 W por \emph{batería} a sesenta minutos, $25\,^{\circ}$C y final 1,80 V/celda. Con eficiencia conjunta mínima propia 0,85, incluida conversión/cable/protección, y reserva de dimensionamiento propia 0,90:
\[
P_{AC,60}=80(1.129)(0,85)(0,90)/1000=69,0948\ \mathrm{kW}>60\ \mathrm{kW}.
\]
Es cálculo de selección desde tabla de batería; no curva oficial de autonomía de la UPS ni extrapolación de envejecimiento. La recepción verificará al menos sesenta minutos iniciales a 60 kW por camino, banco preparado según fabricante y banda térmica de proyecto, con grupo inhibido y ensayo externo. Se programan pruebas individuales, seguimiento de bloques e informe semestral; se sustituye preventivamente al llegar a 45 minutos disponibles o ante degradación anómala, manteniendo el mínimo contractual de treinta minutos. Un autotest no reemplaza la descarga controlada; los bancos A/B no se suman \parencites[pp.~1--2]{csbHRL12540WRA260131}[RT-06.07/10, p.~15]{pucv-transversal}{csbManualJunio2025}.

La sustitución también aplica al alcanzar el criterio de capacidad medida del 80\,\% indicado por CSB, o antes ante tendencia anómala; ese porcentaje no se convierte linealmente en minutos de autonomía. El montaje conserva separaciones de 5--10 mm entre bloques, protección frente a contacto y acceso de elevación. El proyecto de ventilación verifica gasificación, descarga y carga según fabricante, incorporando su aire de renovación al balance sensible/latente. La fuga equivalente de 15 m$^3$/h adoptada para la envolvente no constituye caudal acreditado de ventilación de baterías; si la renovación necesaria excede esa hipótesis, se recalculan clima y envolvente antes de ejecución \parencite{csbManualJunio2025}.

El banco tiene 480 V nominales, flotación propuesta 546 V y techo de carga 588 V, con compensación térmica y perfil autorizado. La tabla CSB permite 540--552 V de flotación y 576--600 V de ecualización para cuarenta bloques; no se aplica automáticamente su extremo superior. La descarga conservadora de cálculo termina a 432 V, antes del corte UPS de 400,8 V. La guía y ficha permiten el intervalo de batería seleccionado; la parametrización se recibe por servicio autorizado. Cada bloque pesa 43,90 kg: 1.756 kg por cadena, 3.512 kg por UPS y 7.024 kg en ambos caminos, sin estructura ni barras. El plano dimensional separa bancos de UPS y TI, con solera y fijación propias.

\textbf{Generación, recarga y recuperación.} Se estudia recarga fija de 20 A DC agregados por UPS tanto en red como en grupo. La guía acredita un máximo configurable de 60 A, pero no confirma el mínimo ni los pasos que permitan fijar exactamente 20 A; el ajuste requiere una configuración autorizada. Los cálculos a 20 A conservan carácter de escenario hasta resolver ese ajuste, la disipación por estado y el tratamiento de bancos. A 588 V representa 11,76 kW DC por UPS, principalmente almacenamiento, no calor íntegro de sala. Para calcular el grupo se imponen eficiencia conjunta de proyecto $\geq0,94$, FP rectificador $\geq0,99$ y auxiliares externos $\leq0,1$ kW a FP $\geq0,75$. Son límites de integración por verificar en sus estados; no datos de catálogo a toda carga parcial \parencite[pp.~3--4]{p7-eaton-93t-tecnica}.
\[
P_g\leq\{48,491615+1,44+1,20(23,845457)(0,02)+11,76\}/0,94+0,1
=66,3382\ \mathrm{kW}.
\]
La cota conservadora con arranque es 67,0406 kVA; con bancada, 66,0829 kW/66,7827 kVA. Cada grupo atiende solo su UPS y su banco. La entrada máxima publicada de 108 A a 400 V equivale a 74,8246 kVA del rectificador completo: al usarla no se vuelve a sumar recarga. Conservar 108 A hasta 478 V daría 89,4154 kVA y, con auxiliares, 89,5487 kVA; es una sensibilidad propia, no curva multivoltaje de fabricante. Se exige demanda total acreditada $\leq90$ kW/90 kVA o recálculo previo a compra.

El grupo deberá conservar capacidad efectiva $\geq129$ kW/$\geq90$ kVA después de correcciones ambientales. El mínimo activo cubre el techo de 90 kW al limitar conservadoramente la media al 70\,\% de PRP efectivo: $129(0,70)=90,3$ kW. El 70\,\% es criterio propio; la ficha P200-6 no se presenta como origen de ese porcentaje. Sus 144 kW nominales se publican a 25\,$^{\circ}$C, 100 m y 30\,\% HR. La capacidad del radiador a 50\,$^{\circ}$C no prueba potencia sin derating. Se requiere selección por temperatura de admisión, altura, ventilación y contrapresión de proyecto; también operación a baja carga, escalones y retorno con banco descargado. No se traslada una rampa de 1 A/s de otra UPS. La regulación estacionaria de grupo $\pm0,5\,\%$ y el arranque de motor a caída de 30\,\% no constituyen ensayo del conjunto \parencites[pp.~1/4--5]{p7-fg-p200-ficha}[p.~3]{p7-eaton-93t-tecnica}.

P200-6 publica 41,3 L/h a 100\,\% PRP de 50 Hz: $24(41,3)(1,20)=1.189,44$ L, redondeados a 1.190 L útiles por grupo. Se especifica estanque nuevo de 1.500 L nominales, descontando expansión y volumen no extraíble según el plano; el tanque estándar de 394 L no cubre la obligación \parencite[pp.~2--3]{p7-fg-p200-ficha}.

\textbf{Servicio de reabastecimiento declarado.} La oferta propone Enex Express como prestador principal de entrega de diésel en terreno y Copec Empresas como alternativa de suministro industrial. Sus servicios públicos respaldan el tipo de prestación; la cobertura específica de Bahía Panitao y los plazos siguientes corresponden a condiciones contractuales propuestas que Synaptix debe formalizar antes de la puesta en servicio \parencites{p6-enex-express}{p6-copec-industrial}.

El servicio cubrirá los dos estanques propios de generación, con atención de pedidos y coordinación de despacho las veinticuatro horas. Synaptix supervisa niveles, activa el pedido, confirma despacho y coordina acceso y recepción con el CLIENTE; un operador competente ejecuta descarga y control de calidad. Cada estanque tendrá al menos 1.190 L útiles. Con consumo de cálculo de 41,3 L/h, esa reserva representa $1.190/41,3=28,81$ horas; sostiene veinticuatro horas sin depender de una entrega durante ese intervalo. El estanque estándar de 394 L del grupo no satisface esta autonomía \parencite[pp.~2--3]{p7-fg-p200-ficha}.

La reposición se activa al llegar a 450 L útiles o cuando la proyección de demanda anticipa ese umbral. El plazo propuesto es de ocho horas desde el pedido, con confirmación en una hora. Si no se confirma, Synaptix activa el prestador alternativo al término de esa hora y conserva el mismo vencimiento de ocho horas contado desde el pedido inicial. A máxima demanda, $450-8(41,3)=119,6$ L proporcionan 2,90 horas adicionales después del vencimiento; la entrega debe restituir al menos 1.190 L útiles por estanque. Nivel, pedido, confirmación, volumen recibido y calidad quedan registrados.

La contratación especificará diésel compatible con el motor, volumen por entrega, acceso terrestre a ambos puntos de descarga, vehículo y personal autorizados, contención, seguridad, limpieza, tratamiento de agua y conciliación de volumen. Antes de una marejada anunciada se completa la reserva y se confirma la ruta alternativa. Si se prevé un bloqueo de acceso superior a la reserva efectiva, se amplía el almacenamiento autorizado o se prepara una provisión segura previa; compartir combustible del surtidor comercial no es la contingencia prevista. El servicio contratado cubre la vigencia completa de operación y se formaliza antes de habilitación, con acceso ordinario y contingencia. La contratación y la prueba de despacho deben acreditar los plazos y la factibilidad local antes de habilitación, sin atribuir al prestador un acuerdo ya firmado.


La energía almacenada se estima con la tensión real del perfil, no con 480 V como techo del cargador. La capacidad de cargador y el cociente energía/potencia no acreditan un plazo de recuperación. Después de descarga, Synaptix conserva generación/reabastecimiento y camino sano, identifica la rama como reserva en recuperación y solo restituye su disponibilidad mediante la relación tensión/corriente/energía y autonomía validada en aceptación. La prueba descarga--recarga--nuevo corte documenta esa relación y la contingencia, sin prometer treinta minutos disponibles por el simple retorno de AC. Los cortes sucesivos no se presentan como autonomía ilimitada; no se añade un plazo numérico de recarga que las bases no exigen \parencites[RT-06.07/08, p.~15]{pucv-transversal}{csbManualJunio2025}{p7-eaton-93t-guia}.

\textbf{Selección climática y estados térmicos.} Se estudian dos circuitos 022/YKA2 con control Evolution+, opciones P051/P161 de humedad, calentador y humidificador 4301. La ficha publica 19,3 kW sensibles brutos menos 0,70 kW de ventilador: 18,60 kW netos a retorno 27\,$^{\circ}$C/47\,\% HR, exterior 35\,$^{\circ}$C, 20 Pa y tubería de 5 m. Sus corrientes de equipo completo ya incluyen auxiliares de esa configuración. La propuesta fija entradas TI/UPS/bancos de 23--25\,$^{\circ}$C, consigna 24\,$^{\circ}$C y HR 40--55\,\%, consigna 45\,\%, sin condensación. Un retorno dirigido hasta 27\,$^{\circ}$C no convierte ese punto en capacidad válida para todo el rango \parencites[p.~3]{p6-mext-pi-2025}{p6-mext-databook}{p6-mext-control-me38}.

Los 2,80092 kW de disipación por 93T corresponden a un punto AC/AC de carga lineal, no a un máximo universal de batería o recarga. Usados como referencia para ambas UPS, con 2,710 kW interiores, envolvente $\leq0,500$ kW y fan de reserva 0,700 kW, la cota con postcalefacción completa de 3,900 kW y margen es $1,25(2,710+2\times2,80092+0,500+0,700+3,900)=16,7648$ kW antes de calor adicional de recarga/bancos. El punto neto nominal permite solo $18,600/1,25-13,41184=1,46816$ kW adicionales. Con eficiencia de cargador de proyecto 0,94, los dos cargadores a 20 A y 588 V sumarían $2(11,76/0,94-11,76)=1,50128$ kW de pérdidas antes del calor de baterías. Por tanto, esta envolvente simultánea todavía excede el punto neto; no se declara resuelta mediante N+1 ni bloqueo de calefacción \parencites[p.~4]{p7-eaton-93t-tecnica}[p.~3]{p6-mext-pi-2025}.

La hipótesis inicial de fugas/ocupación limita aporte de vapor a 0,50 kg/h, sin incluir la reposición de aire de los bancos, equivalente a $0,50(2.450)/3.600=0,3403$ kW latentes. Para exterior de cálculo 35\,$^{\circ}$C/80\,\% HR e interior 24\,$^{\circ}$C/45\,\% a 101,325 kPa, fuga equivalente 15 m$^3$/h aporta 0,33766 kg/h; persona técnica ocasional 0,080 y aperturas adicionales 0,050 suman 0,46766 kg/h. Son condiciones de obra y cálculo, no meteorología observada. La selección debe demostrar extracción media $\geq0,50$ kg/h a carga baja y en la banda T/HR, incluyendo parada, mínimos de compresor y recalentamiento. Los parámetros generales ME38 P12.01/02/03 tienen predeterminados de 180/180/360 s. La lógica específica s-MEXT/PAC-IF mantiene un mínimo ON de tres minutos, retención del paso cero con mínimo OFF y retardos de cinco o diez minutos entre cambios de escalón; reducir P12.01 no acredita eliminar esos mínimos. Sus pasos de mando no representan capacidades uniformes ni modulación térmica continua de cero a cien por ciento \parencite[pp.~139--140 y tabla P12.01--03]{p6-mext-control-me38}. Se requieren capacidad sensible/latente por estado, disipación en batería/recarga, agua compatible con 4301, drenaje, R32, tubería, fijación y protección marina. Las curvas de la variante YKA sin sufijo 2 no se trasladan a YKA2. La memoria de estados y ventilación desarrollada a continuación cuantifica la reposición de bancos y sus cargas adicionales. La configuración final debe conciliar capacidad, tratamiento por zona y recarga autorizada; T-12 conserva RT-06.13 en No hasta completar esa selección.

El recinto común TI/climatización se proyecta con 27,5 m$^2$ e incorpora ambas unidades interiores, sin separación estanca entre sus zonas. Así se individualiza el volumen real de instalación frente al mínimo de 21 m$^2$ de la 022 F2; los bancos mantienen su aire separado. La carga refrigerante y las condiciones de ambas unidades se verificarán con el manual de la pareja suministrada. Esta distribución desarrolla el criterio de superficie y circulación, sin sustituir los balances por estado ni autorizar la compra desde el punto nominal \parencite[secc.~2.2.6--2.2.7, p.~21]{p9-mext-manual-2024}.

El humidificador 4301 requiere agua potable dentro de toda la tabla de calidad del Databook; no se especifica agua desmineralizada ni tratamientos con suavizadores, desinfectantes o inhibidores. La recepción verifica agua, suministro y drenaje del equipo concreto. Esas condiciones, junto con ventilación de bancos y protección marina, forman parte del cierre ambiental antes de compra \parencite[p.~42]{p6-mext-databook}.

\textbf{Tratamiento ENV-25.} Se seleccionan cuatro VEAB CV31-90-3ML con TTC25 externo, dos activos sin solape, con 18 kW de calor y 0,100 kW de controles. El manual permite 40 grados C. Las cuatro cámaras limpias de 2,2 por 1,2 m requieren prevención de retorno de gas y coordinación con serpentín, impulsión y reserva de obra antes de liberar instalación. El frío de proceso calculado, las reservas de potencia y los umbrales de protección no equivalen a una selección húmeda completa ni a autonomía térmica recibida.

\textbf{Desarrollo de tratamiento vigente.} ENV-19 selecciona extracción EX 180A-4 y arquitectura hidráulica exterior; calcula por banco frío de 21,0532 kW hasta rocío de consigna y 22,8604 kW hasta 8,5 grados C, más recalentamiento de 3,5371/4,3215 kW. Los 17,5161 kW netos de reposición no son capacidad suficiente de serpentín frío. El modelo de saturación ideal, los mínimos de proceso y los puntos de placa no sustituyen la selección corregida del chiller, serpentín y fuente de calor. 

\textbf{Balance por estado y ventilación de bancos.} La selección de climatización se construye desde los estados físicos de la cadena, conservando por separado el aire de los bancos y el de cómputo/comunicaciones. La extracción de baterías se conduce al exterior y su reposición se incorpora al tratamiento ambiental de esa zona; no se utiliza el retorno TI como vía de dilución. La Figura~\emph{Plano de distribución proyectada del recinto técnico y zonas asociadas} (SD4) identifica las zonas y la Tabla~\ref{tab:4-estados-termicos} delimita las cargas que deben coexistir. Estas condiciones desarrollan RT-06.13 para la sala secundaria; las consignas de 23--25\,$^{\circ}$C y 40--55\,\% HR son decisiones de proyecto, mientras los límites publicados corresponden a cada variante \parencites[secc.~6.1; RT-06.13, pp.~14--15]{pucv-transversal}{csbManualJunio2025}.

Para una UPS en un estado estacionario, con potencias medidas en sus bornes y signos positivos según el sentido indicado, se emplea
\[
Q_{\mathrm{UPS}}=P_{\mathrm{AC,in}}+P_{\mathrm{DC,desc}}
                -P_{\mathrm{AC,out}}-P_{\mathrm{DC,carga}}.
\]
El término incluye rectificador, inversor, cargador y auxiliares internos. Para el banco se define $Q_{\mathrm{banco}}$ como calor neto cedido al entorno, $E_{\mathrm{alm}}$ como energía interna total y $P_{\mathrm{masa}}$ como energía neta que sale con la materia, incluidos gases. Con tiempos en horas, el balance es $Q_{\mathrm{banco}}=P_{\mathrm{DC,carga}}-P_{\mathrm{DC,desc}}-\mathrm{d}E_{\mathrm{alm}}/\mathrm{d}t-P_{\mathrm{masa}}$. Este calor neto distingue transferencia y generación térmica interna. La selección incluye almacenamiento y procesos térmicos; la resistencia aproximada de 2,60 m$\Omega$ por bloque no proporciona por sí sola un máximo de disipación. En maniobras se añade la variación de energía interna de la UPS. Los 2,80092 kW publicados de la 93T y los factores 0,85/0,90 del cálculo de autonomía cumplen funciones distintas: un punto AC/AC y factores propios de dimensionamiento, respectivamente \parencites{p7-eaton-93t-tecnica}{csbHRL12540WRA260131}.

\begin{table}[htbp]
\centering\small
\caption{Estados que gobiernan la selección térmica}\label{tab:4-estados-termicos}
\begin{tabularx}{\linewidth}{L{25mm}X X}
\hline
\EncabezadoTabla{Estado} & \EncabezadoTabla{Concurrencia física} & \EncabezadoTabla{Comprobación de selección} \\ \hline
Red, carga baja & Dos UPS disponibles, banco en flotación y un circuito climático activo. & Pérdidas de espera y flotación; control de HR con mínimos reales del compresor. \\ \hline
Un camino retirado & UPS sana con el parque completo y su clima; equipos retirados según procedimiento de mantenimiento. & Reparto por fase, calor del camino sano, relevo y suministro individual. \\ \hline
Corte y descarga & Inversor alimentado por su banco; cargador sin aporte AC. & Pérdidas DC/AC y calor de descarga del banco en la banda térmica de autonomía. \\ \hline
Retorno y recarga & Parque en servicio; pueden recuperarse ambos bancos a la vez. & Corriente agregada autorizada, calor de cargadores/bancos y renovación de aire, con reserva en recuperación. \\ \hline
\end{tabularx}
\FuenteTabla{Balance de Synaptix: los extremos se combinan cuando el estado admite su simultaneidad; la retirada de una rama no presupone reparto normal de PSU al 50\,\%.}
\end{table}

Si A descarga y B conserva entrada AC, se calcula ese estado mixto por rama: un mismo banco no recibe simultáneamente la recarga y descarga máximas. Una UPS disponible en reserva conserva su consumo real de espera. Las pruebas de 60 kW se ejecutan por camino y con carga resistiva externa, conservando la otra rama para servicio; no incorporan dos cargas resistivas de 60 kW al calor de la sala TI.

\textbf{Cota propia de gasificación.} CSB recomienda mantener hidrógeno por debajo del 1\,\% volumétrico mediante ventilación. Para evaluar la magnitud sin atribuir una capacidad Ah a la HRL12540W, se calcula un escenario conservador de electrólisis de toda la corriente, sin recombinación. Cuarenta bloques de seis celdas dan $n=240$ celdas en serie. Con $I_\Sigma$ como suma de corrientes de las dos cadenas de una UPS, $F=96.485,33212$ C/mol, $R=8,314462618$ J/(mol K), $T_g=323,15$ K y $p_g=95.000$ Pa:
\[
\dot V_{H_2}=\frac{nI_\Sigma\,3.600}{2F}\frac{RT_g}{p_g},
\qquad x_{H_2}\leq\frac{\dot V_{H_2}}{\dot V_{\mathrm{aire}}}.
\]
Los 50\,$^{\circ}$C y 95 kPa son condiciones conservadoras del volumen gaseoso, distintas del exterior de cálculo térmico; no son consignas de operación del banco. El modelo de mezcla uniforme permite estudiar 300 m$^3$/h efectivos por banco para 20 A: 2,5326 m$^3$/h de H$_2$ y 0,8442\,\%. Para el máximo configurable publicado de 60 A, el escenario escala a 900 m$^3$/h, 7,5978 m$^3$/h y el mismo porcentaje. La corriente agregada evita duplicar el efecto de las cadenas paralelas. Son escenarios propios de cribado, no gasificación medida, caudal normativo ni selección final de ventiladores \parencites{csbManualJunio2025}[p.~98]{p7-eaton-93t-guia}{nistCODATA2022}.

La selección de extracción desarrolla captación alta, ingreso de reposición, mezcla sin bolsas, pérdida de carga con filtros sucios, descarga segura y capacidad conservada ante falla de un ventilador. El control local comprueba caudal e hidrógeno, alerta al NOC y coordina la inhibición de carga mediante una interfaz autorizada, manteniendo la función de descarga que sostenga las cargas. La selección debe identificar esa interfaz: la HMI de la UPS no acredita por sí sola un enclavamiento autónomo. La purga continúa después de cesar carga para evacuar gas residual. La recepción verifica concentración y distribución espacial, flujo efectivo y respuesta ante falla. Un porcentaje calculado bajo mezcla ideal no acredita esas prestaciones constructivas.

\textbf{Efecto sobre sensible y humedad.} Para el escenario de 300 m$^3$/h por banco, ambos bancos requieren 600 m$^3$/h de reposición. Con exterior 35\,$^{\circ}$C/80\,\% HR, interior 24\,$^{\circ}$C/45\,\% HR y 101,325 kPa, se obtiene $w_e=0,0289204$, $w_i=0,0083562$ kg/kg de aire seco y $\rho_{da,e}=1,094636$ kg/m$^3$. El aporte de agua y su carga latente adicional son
\[
\dot m_w=600(1,094636)(0,0289204-0,0083562)
        =13,5062\ \mathrm{kg/h},\qquad
Q_l=13,5062(2.450)/3.600=9,1917\ \mathrm{kW}.
\]
La cota sensible de reposición, con densidad de diseño 1,2 kg/m$^3$ y $c_p=1,020$ kJ/(kg K), es $600(1,2)(1,020)(35-24)/3.600=2,2440$ kW. Estos aportes se añaden a transmisión, fugas residuales, ocupación y calor eléctrico según la zona. El escenario de 900 m$^3$/h por banco triplica los aportes: 40,5187 kg/h, 27,5752 kW latentes y 6,7320 kW sensibles. Si se acredita una gasificación menor y otra renovación mediante datos específicos, se rehace el balance con ese caudal.

Para el tratamiento completo se emplea la diferencia de entalpía del aire húmedo, $h=1,006T+w(2.501+1,86T)$ kJ/kg de aire seco: $600\rho_{da,e}(h_e-h_i)/3.600=11,6774$ kW en ambos bancos para el escenario de 20 A, equivalentes a 5,8387 kW por banco, o 35,0322 kW en ambos para 60 A. Esa carga total incluye el componente latente, por lo que se cuenta una sola vez; los aportes sensible/latente anteriores son aproximaciones de descomposición y no sustituyen el balance entálpico.

Por tanto, el techo de vapor de 0,50 kg/h describe únicamente fugas/ocupación de la hipótesis inicial, sin la reposición de baterías. Incluso el escenario de 20 A añade 9,1917 kW latentes, frente a $22,6-19,3=3,3$ kW latentes en el punto nominal publicado de la 022. La separación de aire evita transferir esa humedad al recinto TI, pero requiere tratamiento propio de bancos y su balance eléctrico; no otorga capacidad adicional a la 022. Del mismo modo, los 23,52 kW DC de dos recargas de 20 A/588 V son potencia entregada a bancos: su fracción térmica depende del estado de almacenamiento, y no se presenta toda esa potencia como disipación real \parencite[p.~3]{p6-mext-pi-2025}.

La combinación 022/YKA2 mantiene carácter de referencia candidata. Su autorización exige selección sensible/latente en la banda completa y carga mínima, disipaciones UPS/banco por estado, extracción y tratamiento compatibles, además del mapa eléctrico y la autonomía. Se compararán reducción del caudal con evidencia específica, tratamiento separado y alternativas de climatización con modulación acreditada, conservando el dimensionamiento proporcional de BT 6.1, C07 RT-06.01 y RT-15.01. Cambiar a una unidad mayor por su placa, o fijar un techo de calor sin respaldo de la variante, no completa la selección. T-12 mantiene el compromiso Sí de RT-06.13, con desarrollo documental no completado por estas causas de diseño; las pruebas futuras verificarán la solución seleccionada, sin sustituirla.


\textbf{Perfil anual y PUE estimada.} La frontera comprende TI y auxiliares interiores, pérdidas de las UPS/bancos y su climatización. La energía exportada a puestos/campo se resta del recinto, conservando las pérdidas que su suministro disipa dentro. Cloud, NOC e industria se evalúan en sus balances respectivos. Para el año de 365 días se adopta temporada alta de noventa días como hipótesis de planificación; los servicios críticos siguen activos 24 horas. La Tabla~\ref{tab:4-perfil-anual} fija horas y fracciones del consumo máximo de compra, no utilización CPU ni consumos medidos \parencites[RT-06.11, p.~15]{pucv-transversal}[cap.~15, RT-21.06]{caso07}.

\begin{table}[htbp]
\centering\small
\caption{Perfil eléctrico anual proyectado}\label{tab:4-perfil-anual}
\begin{tabularx}{\linewidth}{X R{19mm} R{18mm} R{18mm} R{30mm}}
\hline
\EncabezadoTabla{Estado} & \EncabezadoTabla{h/año} & \EncabezadoTabla{TI} & \EncabezadoTabla{Puestos} & \EncabezadoTabla{Motor / móvil} \\ \hline
Alta, atención 18 h/día & 1.620 & 0,65 & 0,70 & 0,080 / 0,50 \\ \hline
Alta, noche 6 h/día & 540 & 0,40 & 0,25 & 0,010 / 0,20 \\ \hline
Resto, atención 12 h/día & 3.300 & 0,45 & 0,55 & 0,040 / 0,35 \\ \hline
Resto, noche 12 h/día & 3.300 & 0,30 & 0,15 & 0,005 / 0,10 \\ \hline
\end{tabularx}
\FuenteTabla{Hipótesis de proyecto de Synaptix sobre horarios de \parencite[cap.~15, RT-21.06]{caso07}; suma $90(18+6)+275(12+12)=8.760$ h. La última columna presenta las fracciones de motor y carga móvil, en ese orden.}
\end{table}

TI interior máxima suma $1,400+0,100+0,500+0,200+0,150=2,350$ kW de servidores, NGFW, core, NVR y WAN. Los auxiliares interiores suman 0,360 kW constantes. En cada estado $i$, $T_i=2,350u_{TI,i}$, $A_i=0,360$ y $X_i=4,214222+0,300+0,600+5,920u_{PC,i}+1,200u_{motor,i}+0,800u_{movil,i}$ kW. $X$ es energía exportada; DC/radio/borde permanecen presupuestados a su cota. Ninguna fracción anual rebaja UPS, banco o grupo a plena carga.

Para la estimación se presupuestan pérdida fija $P_0=0,250$ kW/UPS, coeficiente variable $\eta_v=0,960$, ventilación media conjunta $F=0,900$ kW, rendimiento equivalente de compresión $COP=3,000$, envolvente media $Q_e=0,300$ kW y reserva latente $Q_l=0,100$ kW. Humedad/recalentamiento dispone de $E_H=1.700$ kWh/año y ciclos de bancos de $E_R=250$ kWh/año de pérdidas netas adicionales. Son parámetros energéticos propios para presupuesto y sensibilidad; no prestaciones certificadas de la 93T/022 ni una selección psicrométrica. La recepción sustituirá el modelo por curvas y reparto reales sin sumar nuevamente pérdidas incluidas.

La climatización aumenta salida UPS, y las pérdidas UPS aumentan calor. Se resuelve esa realimentación con $N_i=T_i+A_i+X_i$, $k=1/\eta_v-1$, $H=E_H/8.760$ y $R=E_R/8.760$:
\[
C_i=\frac{F+H+[T_i+A_i+2P_0+kN_i+R+Q_e+F+H+Q_l]/COP}{1-k/COP},
\qquad L_i=2P_0+k(N_i+C_i).
\]
El consumo del recinto es $T_i+A_i+L_i+R+C_i$. El ventilador/humedad se cuenta una vez como electricidad y una vez como calor a retirar, que son fenómenos distintos. En un ciclo anual con estado inicial/final de batería equivalente, la energía útil restituida ya alimentó cargas del perfil: sumar toda la recarga bruta la duplicaría. $E_R$ agrega pérdidas netas no incluidas, y los ensayos externos se registran como exportación y pérdidas adicionales propias.

La Tabla~\ref{tab:4-energia-anual} presenta la energía reconstruible desde ese perfil. La fracción eléctrica media de TI es $8.798,40/(2,350\times8.760)=42,74\,\%$.

\begin{table}[htbp]
\centering\small
\caption{Energía anual y sensibilidad del PUE local}\label{tab:4-energia-anual}
\begin{tabularx}{\linewidth}{X R{30mm}}
\hline
\EncabezadoTabla{Partida del presupuesto inicial} & \EncabezadoTabla{kWh/año} \\ \hline
TI interior & 8.798,40 \\ \hline
Auxiliares interiores & 3.153,60 \\ \hline
Pérdidas de conversión/espera UPS & 8.591,70 \\ \hline
Pérdidas adicionales netas de bancos & 250,00 \\ \hline
Climatización total & 20.877,90 \\ \hline
Recinto técnico, numerador PUE & 41.671,60 \\ \hline
Exportación a puestos/campo, separada & 68.250,87 \\ \hline
\end{tabularx}
\FuenteTabla{Cálculo de Synaptix sin redondear subtotales: $PUE=41.671,60/8.798,40=4,7363\approx4,74$. La energía exportada emplea la cota DC completa $3,7028/0,90+0,100$ kW.}
\end{table}

El PUE de proyecto es 4,74; en crecimiento, al aumentar solo DC/PoE convertido a 5,034222 kW, es 4,78. La cifra refleja TI interior pequeña y dos UPS que también atienden carga exterior, cuyas pérdidas permanecen dentro. Es una estimación de energía anual con hipótesis explícitas, distinta de una razón entre máximos instantáneos. La sensibilidad combinada es 3,42--7,12: caso menor con $P_0=0,150$, $\eta_v=0,975$, $COP=4,0$, $F=0,650$, $E_H=800$, $E_R=150$, $Q_e=0,200$ y $Q_l=0,050$; caso mayor con 0,400, 0,940, 2,2, 1,300, 3.000, 400, 0,500 y 0,200, en las mismas unidades. Se mantienen horas y factores de uso; la amplitud no representa una curva de fabricante ni una predicción meteorológica.

El presupuesto conserva el escenario energético de referencia 022, con medias propias de ventilación, humedad y pérdidas; los escenarios de extracción de la Tabla~\ref{tab:4-estados-termicos} y su memoria no constituyen una instalación ya seleccionada. Al seleccionar caudal, ventiladores y tratamiento de bancos se sustituyen $F$, $Q_e$, $Q_l$, $E_H$ y $E_R$, se incorporan sus estados y se resuelve nuevamente la realimentación. Los 4,74 no se atribuyen como PUE garantizada de una alternativa con equipos o caudales distintos.

La recepción habilita medición separada de entrada por fuente, TI, auxiliares, clima y exportación; en operación se integran doce meses comparables y se actualiza la estimación al cambiar parque o uso. El combustible y emisiones tienen metodología propia. Este PUE desarrolla RT-06.11 junto con kW/FP proyectados; no acredita por sí solo la intensidad de carbono de la región cloud exigida en RT-15.04 ni el control de humedad de RT-06.13.

